\frfilename{mt24.tex}
\versiondate{15.11.13}
\copyrightdate{1995}

\def\chaptername{Ruang fungsi}
\def\sectionname{Pendahuluan}

\newchapter{24}

Kekuatan luar biasa teori integrasi Lebesgue barangkali paling jelas
ditunjukkan oleh kemampuannya menyediakan struktur yang relevan bagi
pertanyaan-pertanyaan yang sangat berbeda dari persoalan yang mula-mula
hendak ditanganinya. Dalam bab ini saya menyajikan konstruksi dan sifat-sifat
elementer beberapa ruang dasar dalam analisis fungsional.

Saya tidak merasa perlu membenarkan kajian tentang ruang bernorma di sini;
jika Anda belum pernah menjumpainya, saya berharap pendahuluan ini setidaknya
akan menunjukkan bahwa ruang-ruang tersebut menyediakan landasan bagi suatu
perpaduan yang mengagumkan antara aljabar dan analisis. Bagian-bagian teori
ruang metrik, ruang bernorma, dan topologi umum yang akan kita perlukan
diuraikan secara ringkas dalam \S\S2A2-2A5.  %2A2 2A3 2A4 2A5
`Ruang-ruang fungsi' utama yang dibahas dalam bab ini sebenarnya memadukan
tiga unsur struktural: ruang-ruang tersebut merupakan ruang linear
(berdimensi tak hingga), ruang metrik, dengan konsep-konsep kontinuitas dan
konvergensi yang menyertainya, serta ruang terurut, dengan pengertian
supremum dan infimum yang bersesuaian. Interaksi di antara ketiga jenis
struktur ini menyediakan khazanah gagasan yang tiada habisnya. Lebih jauh,
banyak dari gagasan tersebut dapat diterapkan secara langsung pada beragam
masalah dalam matematika yang kurang lebih bersifat terapan, khususnya pada
persamaan diferensial dan integral, tetapi secara lebih umum pada setiap
sistem dengan tak hingga banyaknya derajat kebebasan.

Saya menyusun bab ini dengan bagian-bagian mengenai $L^0$ (ruang kelas
ekuivalensi semua fungsi terukur bernilai riil, tempat semua ruang lain dalam
bab ini dibenamkan), $L^1$ (kelas ekuivalensi fungsi-fungsi terintegralkan),
$L^{\infty}$ (kelas ekuivalensi fungsi-fungsi terukur terbatas), dan $L^p$
(kelas ekuivalensi fungsi-fungsi yang pangkat ke-$p$-nya terintegralkan).
Meskipun analisis fungsional biasa jauh lebih banyak memberi perhatian pada
ruang Banach $L^p$ untuk $1\le p\le\infty$ daripada pada $L^0$, dari sudut
pandang khusus buku ini ruang $L^0$ sekurang-kurangnya sama penting dan sama
menariknya dengan ruang-ruang yang lain. Setelah keempat bagian tersebut,
saya kembali mengkaji topologi baku pada $L^0$, yaitu topologi `konvergensi
dalam ukuran' (\S245), kemudian dua bagian yang saling terkait mengenai
keterintegralan seragam dan kekompakan lemah dalam $L^1$ (\S\S246-247).

Ada satu hal teknis di sini yang tidak boleh sekali-kali luput dari perhatian.
Walaupun sudah lazim dan wajar untuk menyebut $L^1$, $L^2$, dan yang lainnya
sebagai `ruang fungsi', unsur-unsurnya sebenarnya bukan fungsi, melainkan
kelas-kelas ekuivalensi fungsi. Seperti terlihat dari bahasa dalam paragraf
sebelumnya, kebiasaan saya adalah mempertahankan pembedaan itu dengan
saksama; alasan saya diberikan dalam catatan untuk \S241.

\discrpage

